\begin{symbols}{p{0.18\textwidth}p{0.7\textwidth}l} % Include a list of Symbols (a three column table)

\textit{Utente} & Si riferisce a qualsiasi utente del sistema e può essere sia \textit{interno} che \textit{esterno}. & \\
\textit{Template di Cockburn} & \`{E} una rappresentazione completa di un caso d'uso. Descrive in modo dettagliato tutti i passaggi necessari all'esecuzione di un caso d'uso e include sezioni di supporto  come \textit{precondizioni} e \textit{postcondizioni}. \\
\textit{Mock-up} & \`{E} un \textit{artefatto}, una rappresentazione visiva statica utilizzata per descrivere un sottoinsieme delle funzionalità del sistema. Il suo fine è puramente illustrativo o espositivo. \\ 
\textit{Issue} & \`{E} una segnalazione di un problema che è stato riscontrato da un membro del team di sviluppo. In altre parole, è un modo per avvisare il team che qualcosa dovrebbe essere sistemato.  \\
\textit{Feature} & In Informatica, si riferisce a una funzionalità o capacità specifica che il software offre ai propri utilizzatori. Nel contesto del prodotto software preso in esame, è un tipo di segnalazione (\textit{issue}) usata per suggerire una nuova funzionalità. \\
\textit{Bug} & In Italiano anche detto \textit{Baco}, nell'ambito della programmazione Informatica, è un'anomalia che porta a un malfunzionamento di un software, producendo un risultato inatteso o errato.\\
\textit{Stakeholder} & Gli \textit{stakeholder} (soggetti interessati o persone coinvolte) sono tutte quelle persone che hanno un interesse nei confronti di un'organizzazione e che influenzano e/o sono influenzate dalle sue attività. \\
\textit{Caso d'uso} & Si riferisce alla descrizione del comportamento del sistema quando riceve uno stimolo da parte di un attore. In UML, un caso d'uso è graficamente rappresentato da un ovale. \\
\textit{Attore} & Specifica un ruolo (o un insieme di ruoli) assunto da un utente o altra entità che interagisce con il sistema. \\
\textit{Container} & \`{E} una tecnologia che consente di creare un ambiente isolato, simulando un ambiente operativo, in cui le applicazioni possono essere eseguite senza preoccuparsi dei conflitti con le altre. \\
\textit{Linguaggi Object-Orientend} & Sono linguaggi appartenenti al paradigma computazionale \textit{Object-Oriented}. Questo pone al centro delle attività la modellazione di concetti del mondo reale e le loro relazioni che definiscono il problema che si vuole risolvere. \\
\textit{CASE} & Significa  \textit{Computer-Aided Software Engineering} (ingegneria del software assistita dal computer) e comprende una vasta gamma di programmi diversi per aiutare le attività di processo del software. \\
\textit{Back-end} & \`{E} la parte di un sistema informatico che gestisce il funzionamento di un'applicazione o sito web. Comprende tutto ciò che avviene dietro le quinte: elaborazione dei dati, logica del sistema, comunicazione della base di dati. \\
\textit{Front-end} & \`{E} la parte di un'applicazione o di un sito web visibile e accessibile all'utente, ovvero l'interfaccia con cui l'utente interagisce. \\
\textit{Docker} & \`{E} una delle piattaforme software container più utilizzate e diffuse nell'ambito della virtualizzazione a livello di sistema operativo. \\
% $L$ & light field function & --- \\
\end{symbols}
